\documentclass[11pt]{article}
\usepackage[margin=2.4cm]{geometry}
\usepackage{fontspec}
\setmainfont{DejaVu Sans}
\usepackage{juliamono}
\usepackage{fvextra}
\usepackage{xcolor}
\usepackage{xurl}
\usepackage{hyperref}
\usepackage{amssymb}
\newcommand{\checkbox}{$\square$}
\newcommand{\answerchoices}{\checkbox\ YES \quad \checkbox\ NO \quad \checkbox\ CANNOT DETERMINE}
\newcommand{\decisionexplanation}{%
  \par\smallskip\textbf{If NO, give the reason for the mismatch (required). If
  CANNOT DETERMINE, state what context or evidence is missing:}\par
  \rule{\linewidth}{0.4pt}\par\vspace{7pt}\rule{\linewidth}{0.4pt}\par}
\setlength{\parindent}{0pt}\setlength{\parskip}{5pt}
\setlength{\emergencystretch}{2em}
\DefineVerbatimEnvironment{Vsrc}{Verbatim}{fontsize=\small,frame=single,
  breaklines=true,breakanywhere=true}
\DefineVerbatimEnvironment{Vlean}{Verbatim}{fontsize=\scriptsize,frame=single,
  numbers=left,numbersep=5pt,breaklines=true,breakanywhere=true}
\hypersetup{colorlinks=true,urlcolor=blue}

\begin{document}
\begin{center}{\LARGE Proof Verification Document}\\[2pt]{\large proofwiki\_69483}\\[2pt]{\small ProofWiki flagged proofs (newly confirmed)}\end{center}
\begin{center}\fcolorbox{blue!60!black}{blue!4}{\parbox{0.92\linewidth}{\textbf{Make three independent decisions.} Follow the evidence chain: original source $\longrightarrow$ challenged claim $\longrightarrow$ proposed counterexample $\longrightarrow$ Lean theorem. A record is valid only when all three links are correct. Complete the human sign-off before reading the automated review appendix; those earlier results are background, not an answer key.}}\end{center}
\textbf{Use only this document, except for the public source link.} Everything needed for Questions 2 and 3 is included here. For Question 1, open the public source link when one is shown. If there is no public link, the complete stored problem and solution are included here. If the information is still not enough, choose CANNOT DETERMINE.
For each question, choose \textbf{YES}, \textbf{NO}, or \textbf{CANNOT DETERMINE}. Choose CANNOT DETERMINE if the document does not provide enough information for you to decide; do not mark YES if you are uncertain. If you find a specific mismatch, choose NO and explain it in the space below that question. Do not rewrite the claim, add assumptions, or substitute a different counterexample.
\textbf{How copied text is shown.} The problem, quoted passage, and proposed counterexample come from the stored record. The Lean theorem statement and Lean files are copied exactly from the certificate. Some long lines continue on the next printed line; this is only page formatting and does not change the copied text. Line numbers are added only to help you find passages. If the stored record contains a character that cannot be printed, the document shows \texttt{[U+XXXX CONTROL]} in its place.
\begin{description}
\item[Challenged claim.] The exact statement from the source that this record says is false.
\item[Proposed counterexample.] A concrete case that should satisfy the claim's assumptions but make its conclusion false.
\item[Lean certificate.] A Lean file proving a formal theorem. Successful compilation proves the theorem as written, not that it is the right translation of the source.
\end{description}
\section*{Question 1 --- Does this proof verification document accurately represent the source?}
\textbf{What to decide.} Compare the challenged passage below with the original source---either the public link or the complete source text in this document. Answer YES only if they make the same claim and include the same assumptions or conditions. If this document changes the meaning, leaves out a condition, or points to a different passage, answer NO.
\subsection*{Original source}
\textbf{Source record:}
\begin{Vsrc}
qedbench_certificate_availability_triage/data/records_proofwiki.jsonl
\end{Vsrc}
\textbf{Public source:}\par
\url{https://proofwiki.org/wiki/Summation_from_k_to_m_of_r_Choose_k_by_s_Choose_n-k_by_nr-(r\%2Bs)k}\par
\subsection*{Problem or question from the source (verbatim)}
\begin{Vsrc}
Title: Summation from k to m of r Choose k by s Choose n-k by nr-(r+s)k

:$\ds \sum_{k \mathop = 0}^m \dbinom r k \dbinom s {n - k} \paren {n r - \paren {r + s} k} = \paren {m + 1} \paren {n - m} \dbinom r {m + 1} \dbinom s {n - m}$
\end{Vsrc}
\subsection*{Challenged passage recorded from the source (stored text)}
\begin{Vsrc}
\paren {j + 2} \paren {r - j + 2} \dbinom r {j + 2} \paren {\paren {j + 1} \paren {n - j} \dbinom s {n - j} + \dbinom s {n - j - 1} \paren {n r - \paren {r + s} \paren {j + 1} } }
\end{Vsrc}
\subsection*{Surrounding source lines (stored text with added line labels)}
The labels refer to lines in the solution text stored with the source record; the source-line content is not transliterated.
\begin{Vsrc}
0043 |       | o = 
0044 |       | c = 
0045 | }}
0046 | {{eqn | r = \sum_{k \mathop = 0}^j \dbinom r k \dbinom s {n - k} \paren {n r - \paren {r + s} k} + \dbinom r {j + 1} \dbinom s {n - j - 1} \paren {n r - \paren {r + s} \paren {j + 1} }
0047 |       | c = 
0048 | }}
0049 | {{eqn | r = \paren {j + 1} \paren {n - j} \dbinom r {j + 1} \dbinom s {n - j} + \dbinom r {j + 1} \dbinom s {n - j - 1} \paren {n r - \paren {r + s} \paren {j + 1} }
0050 |       | c = 
0051 | }}
0052 | {{eqn | r = \dbinom r {j + 1} \paren {\paren {j + 1} \paren {n - j} \dbinom s {n - j} + \dbinom s {n - j - 1} \paren {n r - \paren {r + s} \paren {j + 1} } }
0053 |       | c = factorising
0054 | }}
0055 | {{eqn | r = \paren {j + 2} \paren {r - j + 2} \dbinom r {j + 2} \paren {\paren {j + 1} \paren {n - j} \dbinom s {n - j} + \dbinom s {n - j - 1} \paren {n r - \paren {r + s} \paren {j + 1} } }
0056 |       | c = {{Corollary|Factors of Binomial Coefficient|2}}
0057 | }}
0058 | {{eqn | r = \paren {j + 2} \paren {r - j + 2} \dbinom r {j + 2} \paren {\dfrac {\paren {j + 1} \paren {n - j} } {\paren {n - j} \paren {s - \paren {n - j} } } \dbinom s {n - j - 1} + \dbinom s {n - j - 1} \paren {n r - \paren {r + s} \paren {j + 1} } }
0059 |       | c = {{Corollary|Factors of Binomial Coefficient|2}}
0060 | }}
0061 | {{eqn | r = \paren {j + 2} \paren {r - j + 2} \dbinom r {j + 2} \dbinom s {n - j - 1} \paren {\dfrac {j + 1} {s - n + j} + \paren {n r - \paren {r + s} \paren {j + 1} } }
0062 |       | c = simplifying
0063 | }}
0064 | {{eqn | r = \paren {j + 2} \paren {r - j + 2} \dbinom r {j + 2} \dbinom s {n - j - 1} \paren {\dfrac {j + 1} {s - n + j} + \paren {n r - r j - s j - r - s} }
0065 |       | c = simplifying
0066 | }}
0067 | {{end-eqn}}
\end{Vsrc}
\textbf{A YES means:} the proof verification document faithfully represents both the source context and the challenged claim, without dropping a condition or changing its meaning.
\section*{Question 2 --- Does the proposed counterexample disprove that claim?}
\textbf{Claim to test (repeated for convenience).}
\begin{Vsrc}
\paren {j + 2} \paren {r - j + 2} \dbinom r {j + 2} \paren {\paren {j + 1} \paren {n - j} \dbinom s {n - j} + \dbinom s {n - j - 1} \paren {n r - \paren {r + s} \paren {j + 1} } }
\end{Vsrc}
\subsection*{How the record says the argument uses this claim}
This is the record's description, not an established fact; check it against the source.
\begin{Vsrc}
In the induction step, the transition from line 0052 to line 0055 (justified by 'Corollary: Factors of Binomial Coefficient 2') asserts the identity C(r, j+1) = (j+2)(r-j+2) C(r, j+2), i.e. it replaces the factor \dbinom r {j+1} by \paren{j+2}\paren{r-j+2}\dbinom r {j+2} while leaving the bracketed sum unchanged. This local equality is false.
\end{Vsrc}
\subsection*{Proposed counterexample (exact values or construction)}
\begin{Vsrc}
Take r = 2, j = 0 (so j+1 = 1, j+2 = 2, both binomials in range and strictly positive). The step from line 0052 to 0055 asserts C(r, j+1) = (j+2)(r-j+2) C(r, j+2). Here the left side is C(2,1) = 2, while the right side is (j+2)(r-j+2)C(2,2) = 2 * 4 * 1 = 8. So the claim says X = Y with X = 2 and Y = 8: 2 != 8, a factor-of-4 discrepancy. Equivalent smaller/other witnesses: r = 3, j = 1 gives C(3,2) = 3 vs (3)(4)C(3,3) = 12; r = 4, j = 1 gives C(4,2) = 6 vs (3)(5)C(4,3) = 60. In the generalized (rational upper index) setting r = -40, j = -1 gives C(r,0) = 1 vs (1)(-37)... = 1480, so the failure is not an artifact of integrality.
\end{Vsrc}
\subsection*{Why the record says the counterexample works}
Check this explanation independently; it is evidence to assess, not an answer key.
\begin{Vsrc}
The flagged assertion is that the local equality C(r, j+1) = (j+2)(r-j+2) C(r, j+2), implicitly used when line 0052 becomes line 0055, is false. Because the bracketed sum ((j+1)(n-j)C(s,n-j) + C(s,n-j-1)(nr-(r+s)(j+1))) is carried across unchanged, the two lines agree only if that binomial replacement is value-preserving, so a single pair (r, j) violating it refutes the step. The script evaluates both sides in exact integer arithmetic (math.comb with the k<0 or k>n => 0 convention) over all 14883 pairs with 0 <= r <= 120 and -2 <= j <= 120, and in exact Fraction arithmetic over 4860 generalized-binomial triples. Result: the equality holds on 7502 pairs, but every one of those is degenerate (both sides zero because a binomial is out of range). Among the 7260 non-degenerate pairs (0 <= j+1 and j+2 <= r, both binomials strictly positive), the identity fails on all 7260 — it holds on NONE. The structural reason, consistent with the exhaustive scan: the genuine "Factors of Binomial Coefficient" corollary is (j+2) C(r, j+2) = (r-j-1) C(r, j+1), i.e. C(r, j+1) = [(j+2)/(r-j-1)] C(r, j+2). The proof's multiplier (j+2)(r-j+2) equals the true multiplier (j+2)/(r-j-1) only if (r-j+2)(r-j-1) = 1, which is impossible for integers since with d = r-j >= 2 one has (d+2)(d-1) >= 4. So the proof has turned a division by (r-j-1) into a multiplication by (r-j+2); the step is invalid, not merely unproved. No floats were used anywhere; the search is deterministic and stdlib-only.
\end{Vsrc}
\subsection*{How to decide}
\begin{enumerate}
\item List every assumption that applies to the challenged claim.
\item Check that the proposed example satisfies every one of those assumptions.
\item Recompute the claimed outcome and the actual outcome using exact reasoning.
\item Confirm that the conflict is with the challenged claim itself, not an unrelated sentence.
\end{enumerate}
\textbf{A YES means:} the example is fully specified, satisfies all relevant assumptions, and makes the challenged claim false. If a necessary value, condition, or calculation is missing, choose CANNOT DETERMINE.
\section*{Question 3 --- Does the final Lean theorem statement verify this same counterexample?}
\textbf{What to decide.} A Lean file can compile while proving the wrong theorem. Compare these three links independently:
\begin{enumerate}
\item The Lean definitions represent the same mathematical objects, assumptions, and operations as the challenged claim.
\item The Lean construction uses the same concrete values or example shown in Question 2.
\item The final theorem statement asserts that the challenged claim fails on that example, not merely a weaker or nearby statement.
\end{enumerate}
\subsection*{Selected final Lean theorem statement (verbatim UTF-8 excerpt)}
This is an exact contiguous slice of the certificate file, from the theorem keyword through the original \texttt{:=} token where its proof begins. It is shown first so the formal target is not buried in the complete file.
\begin{Vlean}
theorem refutes_line_0052_to_0055 :
    1 <= j /\
    j + 1 <= r /\
    j + 2 <= r /\
    choose r (j + 1) > 0 /\
    choose r (j + 2) > 0 /\
    choose r (j + 1) = 3 /\
    choose r (j + 2) = 1 /\
    unchangedBracket = 3 /\
    line0052 = 9 /\
    line0055 = 36 /\
    line0052 != line0055 /\
    Not assertedLocalEquality /\
    Not assertedLineTransition :=
\end{Vlean}
\subsection*{Project-local imported Lean files (complete, verbatim UTF-8)}
This section follows the certificate's import commands recursively and prints every repository-owned Lean source file on which the final theorem statement depends. Each file is copied in full, including its definitions and controls. External Mathlib, Std, and Lean Init source is not duplicated.
No project-local imported Lean files were found. External Mathlib, Std, and Lean Init source is not duplicated in this proof verification document.
\subsection*{Complete counterexample Lean file (verbatim UTF-8)}
The complete Lean file is included below so you can check every definition and supporting result used by the final theorem statement. Use it together with the project-local imported Lean files above.
\textbf{Lean certificate file:}
\begin{Vsrc}
partB_feasibility/diverse_certificates/proofwiki_69483.lean
\end{Vsrc}
\begin{Vlean}
/-!
Record: proofwiki_69483
Source: ProofWiki, "Summation from k to m of r Choose k by s Choose n-k by
nr-(r+s)k"
https://proofwiki.org/wiki/Summation_from_k_to_m_of_r_Choose_k_by_s_Choose_n-k_by_nr-(r%2Bs)k

Exact target: in the induction step from line 0052 to line 0055, the proof
replaces `C(r,j+1)` by `(j+2)(r-j+2)C(r,j+2)` while leaving the bracket
`(j+1)(n-j)C(s,n-j) + C(s,n-j-1)(nr-(r+s)(j+1))` unchanged.  Thus the local
claim is `C(r,j+1) = (j+2)(r-j+2)C(r,j+2)`; this certificate also evaluates
the complete adjacent expressions.

Witness: `r = 3`, `j = 1`, `n = 3`, and `s = 3`.  Both binomial indices are
in range, both coefficients are positive, and the source-context condition
`j >= 1` holds.  The unchanged bracket is 3, so the local replacement is
`3 = 12`, and the complete lines are 9 and 36.  The earlier Terra witness
`r = 2, j = 0` showed the same arithmetic defect but missed `j >= 1`; the
semantic audit's corrected in-context witness is used here.

Scope: this refutes only the material local proof step, not the page's main
summation claim.  Terra's corrected relation and negative control is
`(j+2)C(r,j+2) = (r-j-1)C(r,j+1)`, which evaluates here to `3 = 3`.

Modelling judgments: natural binomial coefficients are defined locally by
Pascal recursion.  Parameters and indices are natural numbers, while the
bracket and complete line expressions use integers so that its subtraction
has the source's signed meaning.  Natural subtraction is used only for the
nonnegative indices and factors certified in range below.  All witness facts
are closed, exactly computable propositions checked by kernel `decide`; there
are no imports.
-/

namespace ProofWiki69483

def choose : Nat -> Nat -> Nat
  | _, 0 => 1
  | 0, _ + 1 => 0
  | n + 1, k + 1 => choose n k + choose n (k + 1)

def r : Nat := 3
def j : Nat := 1
def n : Nat := 3
def s : Nat := 3

def unchangedBracket : Int :=
  ((j + 1 : Nat) : Int) * ((n - j : Nat) : Int) *
    (choose s (n - j) : Int) +
  (choose s (n - j - 1) : Int) *
    ((n : Int) * (r : Int) -
      ((r + s : Nat) : Int) * ((j + 1 : Nat) : Int))

def line0052 : Int :=
  (choose r (j + 1) : Int) * unchangedBracket

def line0055 : Int :=
  ((j + 2 : Nat) : Int) * ((r - j + 2 : Nat) : Int) *
    (choose r (j + 2) : Int) * unchangedBracket

abbrev assertedLocalEquality : Prop :=
  (choose r (j + 1) : Int) =
    ((j + 2 : Nat) : Int) * ((r - j + 2 : Nat) : Int) *
      (choose r (j + 2) : Int)

abbrev assertedLineTransition : Prop :=
  line0052 = line0055

theorem refutes_line_0052_to_0055 :
    1 <= j /\
    j + 1 <= r /\
    j + 2 <= r /\
    choose r (j + 1) > 0 /\
    choose r (j + 2) > 0 /\
    choose r (j + 1) = 3 /\
    choose r (j + 2) = 1 /\
    unchangedBracket = 3 /\
    line0052 = 9 /\
    line0055 = 36 /\
    line0052 != line0055 /\
    Not assertedLocalEquality /\
    Not assertedLineTransition := by
  decide

/- Negative control: the corrected nearby binomial identity holds at the
same non-degenerate witness, with both sides equal to 3.  This would fail if
the encoding made all relevant binomial equalities uniformly false. -/
theorem negative_control_corrected_binomial_relation_holds :
    ((j + 2 : Nat) : Int) * (choose r (j + 2) : Int) = 3 /\
    (r - j - 1 : Nat) * choose r (j + 1) = 3 /\
    (j + 2) * choose r (j + 2) =
      (r - j - 1) * choose r (j + 1) := by
  decide

end ProofWiki69483

\end{Vlean}
\subsection*{Proof-safety information}
The automated scan found no unfinished proofs, new axiom declarations, unsafe declarations, external code hooks, or uses of Lean's compiled decision procedure. This scan concerns proof construction only; it does not establish that the theorem matches the source claim.
\textit{Technical scan counts:} native\_decide=0, sorry=0, axiom\_decl=0, unsafe=0, extern=0. \textit{Technical tier:} KERNEL-ONLY.
\textbf{A YES means:} Lean represents the same objects and assumptions, uses the same proposed counterexample, and the final theorem statement asserts the same contradiction. Successful compilation by itself is not enough.
\clearpage
\section*{Human sign-off}
Complete this block from your own review of Questions 1--3. Do not use the automated appendix below as an answer key.\par
\textbf{Reviewer name:} \rule{9cm}{0.4pt}\par
\textbf{Date:} \rule{4cm}{0.4pt}\par
\medskip\textbf{1. Does the proof verification document accurately represent the source and challenged claim?}\par
\answerchoices\par
\decisionexplanation
\medskip\textbf{2. Does the proposed counterexample satisfy the assumptions and make that claim false?}\par
\answerchoices\par
\decisionexplanation
\medskip\textbf{3. Does the final Lean theorem statement represent and verify that same counterexample?}\par
\answerchoices\par
\decisionexplanation
\medskip\textbf{Overall decision:}\par
\checkbox\ VALIDATED (all three answers are YES)\par
\checkbox\ REJECTED (at least one answer is NO)\par
\checkbox\ NEEDS ANOTHER REVIEWER (at least one CANNOT DETERMINE and no NO)\par
\medskip\textbf{Additional comments.} Use this space for any extra source, counterexample, calculation, or Lean details that support your decisions.\par\rule{\linewidth}{0.4pt}\par\vspace{7pt}\rule{\linewidth}{0.4pt}
\clearpage
\section*{Appendix --- Prior automated review (background only)}
\textbf{These results are not an answer key.} They come from the earlier Opus/Fable verification campaign and are printed only so an auditor can see what was previously checked. The human decisions above must be made independently. Raw status codes are retained in small type for traceability.
\textbf{Overall automated disposition:} Cleared by automated screening; an independent search also found a counterexample.
\par\textit{Audit code:} \nolinkurl{PASS_CORROBORATED}
\subsection*{Independent search with no access to the stored example}
\textbf{Result:} An independent search found a counterexample without seeing the stored example.
\par\textit{Audit code:} \nolinkurl{WITNESS_FOUND}
\par\textbf{Cases checked:}
\begin{Vsrc}
r,s,n in 0..14 and j in -1..14; 54000 tuples swept exhaustively. Binomials with nonnegative-integer upper index, C(a,b)=0 for b<0 or b>a. Sanity sweeps on the same grid: ambient identity failures 0/ (15^3 x 15 (m in 0..14)); V!= T failures 0/54000.
\end{Vsrc}
\subsection*{Independent re-check of the proposed counterexample}
\textbf{Result:} A separate checker recomputed the proposed counterexample and confirmed it.
\par\textit{Audit code:} \nolinkurl{CONFIRMED}
\subsection*{Independent review of the source-to-Lean correspondence}
\textbf{Result:} A separate reviewer judged that Lean proves the same failure using the same example.
\par\textit{Audit code:} \nolinkurl{MATCHES}
\par\textbf{Theorem summary from that reviewer:}
\begin{Vsrc}
In namespace ProofWiki69483 the file defines `choose : Nat -> Nat -> Nat` by the Pascal recursion (`choose _ 0 = 1`, `choose 0 (k+1) = 0`, `choose (n+1) (k+1) = choose n k + choose n (k+1)`, so the out-of-range convention C(n,k)=0 for k>n is built in), fixes the closed constants `r = 3`, `j = 1`, `n = 3`, `s = 3`, and defines over Int: `unchangedBracket = (j+1)*(n-j)*choose s (n-j) + choose s (n-j-1) * (n*r - (r+s)*(j+1))`, `line0052 = choose r (j+1) * unchangedBracket`, and `line0055 =. (Shortened here; the complete note remains in the machine-check JSON.)
\end{Vsrc}
\par\textbf{Difference noted by that reviewer:}
\begin{Vsrc}
None that break the correspondence; all three links check out, and I re-evaluated every `decide` obligation independently in Python with exact integer arithmetic. (1) Definitions: `choose` by Pascal recursion is the ordinary binomial coefficient with C(n,k)=0 for k>n (spot-checked: choose 3 2 = 3, choose 3 3 = 1), matching the record's stated math.comb-with-zero-convention. (Shortened here; the complete note remains in the machine-check JSON.)
\end{Vsrc}
\subsection*{Follow-up review}
No follow-up review was needed or recorded.
\textit{This appendix is only background. You do not need any separate machine-output files to answer the three questions.}
\end{document}